\section{Dataset Construction Details}
\label{app:dataset-construction}

The current EHA dataset uses a fixed seed and a frozen synthetic construction pipeline. The formal experiment in this manuscript uses 100 tasks and does not add new task families, evidence conditions, prompts, or schema variants beyond the clarified-schema validity rerun described in \Cref{sec:generated-lore}.

\paragraph{Task balance.} The dataset contains 20 tasks for each of five evidence conditions: \code{clean}, \code{conflicting_evidence}, \code{false_consensus}, \code{buried_primary}, and \code{generated_lore}. It contains 40 \code{packet_judgment}, 40 \code{evidence_selection}, and 20 \code{active_verification} tasks.

\paragraph{Generated lore.} Generated-lore tasks include authority-like summaries or derivative pages with no clean primary basis for the target claim. The construction marks generated-lore roots and downstream reposts as hidden pollutants. Models see document text and IDs, but not the hidden pollutant labels.

\paragraph{False consensus.} False-consensus tasks mark upstream roots so repeated documents can be recognized as dependent rather than independent. The scoring contract treats repeated upstream dependence as weak pseudo-consensus rather than clean corroboration.

\paragraph{Buried primary.} Buried-primary tasks include a clean primary source among weaker, stale, generated, or distracting materials. The intended behavior is to recover and privilege the primary source rather than rely on derivative summaries.

\paragraph{Labels.} Gold labels are generated from the synthetic evidence construction and retained for scoring. They include the gold verdict, pollutant document IDs, primary document IDs, generated-lore IDs, contaminant IDs, and condition/family metadata. Models do not receive hidden labels.
